\makeatletter
\hbadness = 99999
\binoppenalty=\maxdimen
\relpenalty=\maxdimen

% \zzcommand{commandName}
\def\zzcommand#1{\let#1\undefined\newcommand#1}

% \DeclareCaptionFormat{capformat}{\fontsize{11pt}{13pt}\selectfont#1#2#3}
% \captionsetup[figure]{textfont=it, labelfont=bf, format=capformat}
% \zzcommand{\floatpagefraction}{.8}
% \zzcommand{\topfraction}{.9}
% \zzcommand{\bottomfraction}{.8}
% \zzcommand{\textfraction}{.1}
% \setcounter{totalnumber}{5}
% \setcounter{topnumber}{5}
% \setcounter{bottomnumber}{5}

% \setcounter{secnumdepth}{3}


\let\Oldincludegraphics\includegraphics
\RenewDocumentCommand{\includegraphics}{O{} m}{%
	\begin{adjustbox}{max width=\linewidth, 
        % max height=\textheight-5\baselineskip
        }
	    \Oldincludegraphics[#1]{#2}
	\end{adjustbox}
}

\let\origfigure\figure
\let\endorigfigure\endfigure

% \renewenvironment{figure}[1]{ %[tbp]{%
%     % \captionsetup{width=\linewidth}
%     \origfigure[#1]%
%     % \centering
% }{%
%     \endorigfigure
% }

\SetKwRepeat{DoWhile}{do}{while}

% Macros to hold current titles
% \currentsection
\zzcommand{\currentsection}{}
% \currentsubsection
\zzcommand{\currentsubsection}{}
% \currentsubsubsection
\zzcommand{\currentsubsubsection}{}


% Save original commands
\let\oldsection\section
\let\oldsubsection\subsection
\let\oldsubsubsection\subsubsection


\gdef\currentsection{}
\gdef\currentsubsection{}
\gdef\currentsubsubsection{}
\zzcommand{\section}[1]{
  \gdef\currentsection{}
  \gdef\currentsubsection{}
  \gdef\currentsubsubsection{}
  \ifthenelse{\equal{#1}{*}}%
  {\@sectionstar}%
  {\@sectionnostar{#1}}%
  % \@ifstar{\@sectionstar}{\@sectionnostar}
}
\zzcommand{\@sectionnostar}[1]{%
  \gdef\currentsection{#1}%
  \oldsection{#1}%
}
\zzcommand{\@sectionstar}[1]{%
  \gdef\currentsection{#1}
  \oldsection*{#1}%
}
\zzcommand{\subsection}[1]{
  \gdef\currentsubsection{}
  \gdef\currentsubsubsection{}
  \ifthenelse{\equal{#1}{*}}%
  {\@subsectionstar}
  {\@subsectionnostar{#1}}
}
\zzcommand{\@subsectionnostar}[1]{%
  \gdef\currentsubsection{#1}%
  \oldsubsection{#1}%
}
\zzcommand{\@subsectionstar}[1]{%
  \gdef\currentsubsection{#1}%
  \oldsubsection*{#1}%
}
\zzcommand{\subsubsection}[1]{
  \gdef\currentsubsubsection{}
  \ifthenelse{\equal{#1}{*}}%
  {\@subsubsectionstar}
  {\@subsubsectionnostar{#1}}
}
\zzcommand{\@subsubsectionnostar}[1]{%
  \gdef\currentsubsubsection{#1}%
  \oldsubsubsection{#1}%
}
\zzcommand{\@subsubsectionstar}[1]{%
  \gdef\currentsubsubsection{#1}%
  \oldsubsubsection*{#1}%
}




% \LetLtxMacro{\Oldincludegraphics}{\includegraphics}
% \RenewDocumentCommand{\includegraphics}{O{} m}{%
% 	\begin{adjustbox}{max width=\linewidth, max height=\textheight-5\baselineskip}
% 	    \Oldincludegraphics[#1]{#2}
% 	\end{adjustbox}
% }
% \inset{inset_size}{path_to_main_image}{path_to_inset_image}
\zzcommand{\inset}[3][0.2\linewidth]{
    \noindent\stackinset{r}{0.1cm}{b}{0.1cm}{
        \colorbox{white}{
            \includegraphics[width=#1]{#3}
        }
    }
	{\includegraphics{#2}}
}

\let\origfigure\figure
\let\endorigfigure\endfigure

\renewenvironment{figure}[1][tbp]{%
    \captionsetup{width=\linewidth}
    \origfigure[#1]%
    \centering
}{%
    \endorigfigure
}

% \teaser{/* Inside commands of the figure* environment */}
\zzcommand{\teaser}[1]{
  \begin{figure*}[ht]
    \centering
        #1
    \end{figure*}
}

% \minitoc
\def\minitoc{
	\etocsettocstyle{\section*{\contentsname}}{}
	\localtableofcontents
	% \noindent\textbf{\hyperlink{tocpage}{$\uparrow$ Back to summary}}
}
% \@addtoreset{chapter}{part}%

\sisetup{product-units=repeat}

\emergencystretch=1em

\zzcommand{\cellalign}{cl}
\newcolumntype{H}{>{\setbox0=\hbox\bgroup}c<{\egroup}@{}}





\zzcommand{\citeName}{\citeauthor}

% % "New" citation style
% \DeclareCiteCommand{\cite}[]
%   {\usebibmacro{prenote}\bibopenbracket}%
%   {\usebibmacro{citeindex}\usebibmacro{countcite}{\thefield{entrykey}}%
%   \printtext[bibhyperref]{\printfield{labelalpha}}}%  \printtext[bibhyperref]{\usebibmacro{cite}}}%
%   {\multicitedelim}%
%   {\bibclosebracket\usebibmacro{postnote}}

% \newbibmacro*{citefullbody}{%
%   \printnames{labelname}%
%   \setunit{\addcomma\space}%
%   \textit{\printfield[]{title}}~(\printfield{year})%
% }

% \DeclareCiteCommand{\citep}
% {\usebibmacro{prenote}}%
% {\usebibmacro{citeindex}\usebibmacro{countcite}{\thefield{entrykey}}%
% \printtext[bibhyperref]{\usebibmacro{citefullbody}}}%
% {\multicitedelim}%
% {\usebibmacro{postnote}}


% % Only print prefix (von/de/van…) + family; enable "et al."
% \DeclareNameFormat{onlyfamily}{%
%   \nameparts{#1}%
%   % ← This is the crucial line: allow biblatex to insert name list delimiters
%   \usebibmacro{name:delim}{\namepartfamily}%
%   % print prefix (if any) and family
%   \ifdefvoid{\namepartprefix}{}{\namepartprefix\space}%
%   \namepartfamily%
%   % allow truncation to print "et al."
%   \usebibmacro{name:andothers}%
% }

% % Use that format for the labelname list (author or editor fallback)
% \DeclareNameAlias{labelname}{onlyfamily}

% % Delimiters BETWEEN names (global)
% \DeclareDelimFormat{finalnamedelim}{\addspace\bibstring{and}\space}
% \DeclareDelimFormat{multinamedelim}{\addcomma\space}
% \DeclareDelimFormat{andothersdelim}{\addspace}

% % Plain-text cite command
% \DeclareCiteCommand{\citeName}
%   {%
%     \usebibmacro{prenote}%
%     \defcounter{minnames}{1}%
%     \defcounter{maxnames}{2}% 1 author → 1; 2 authors → both; ≥3 → "et al."
%   }
%   {%
%     \usebibmacro{citeindex}%
%     \printnames{labelname}%
%   }
%   {\multicitedelim}
%   {\usebibmacro{postnote}}


% \DeclareCiteCommand{\citeProgram}[]
%   {\usebibmacro{prenote}}
%   % {\usebibmacro{citeindex}\usebibmacro{countcite}{\thefield{entrykey}}\printtext[bibhyperref]{\mkbibemph{\printfield{title}},~\printdate}}
%   {\usebibmacro{citeindex}\usebibmacro{countcite}{\thefield{entrykey}}
%   % \printtext[bibhyperref]{\printfield{title}~(\printdate)}}
%   \printtext[bibhyperref]{\printfield{title}}}
%   {\multicitedelim}
%   {\usebibmacro{postnote}}

% % \AtEveryBibitem{\clearfield{url}}
% \AtEveryBibitem{\clearfield{abstract}}
% \AtEveryBibitem{\clearfield{doi}}
% \AtEveryBibitem{\clearfield{note}} % Notes from Mendeley
% \newtotcounter{citenum}
% \AtEveryBibitem{\stepcounter{citenum}}


% \renewbibmacro*{title}{%
% {}
% {\printtext[title]{\mkbibbold{\printfield[title]{title}}}\newunit}}

% \DefineBibliographyStrings{english}{
%   backrefpage = {Cited on page},
%   backrefpages = {Cited on pages},
% }



% \setlength\bibitemsep{1.2\itemsep}


\newcounter{AltTextImageCurrentSide}
\newcounter{ItemizeDepthCounter}
\newlength{\ItemizeLeftMargin}
\newlength{\ItemizeIndentPerLevel}
\newlength{\ItemizeExtraIndent}
\setlength{\ItemizeIndentPerLevel}{1em}

\newenvironment{Itemize}
{%
  \setcounter{AltTextImageCurrentSide}{0}
  \stepcounter{ItemizeDepthCounter}
  \setlength{\ItemizeLeftMargin}{0pt}
  \setlength{\ItemizeExtraIndent}{\ItemizeIndentPerLevel}
  \multiply\ItemizeExtraIndent by \numexpr\value{ItemizeDepthCounter}-1\relax
  \addtolength{\ItemizeLeftMargin}{\ItemizeExtraIndent}
  \begin{itemize}[
    label={}, 
    leftmargin=\ItemizeLeftMargin,
    itemindent=1em,
    itemsep=0em,
    topsep=0em
  ]
}
{%
  \end{itemize}
  \addtocounter{ItemizeDepthCounter}{-1}
  \ifthenelse{\value{ItemizeDepthCounter} > 0}{}{\bigbreak}
}

% Define new command \Itemthat applies italics to the title only
% \Item{} \textbf{text}
\zzcommand{\Item}[1]{%
  \ifblank{#1}%
    {\item $\bullet$ }% Case when #1 is empty
    {\item \textbf{#1}}% Case when #1 is not empty
}


\newcounter{alttextimage@saved@equation}
\newcommand{\ATI@suspend}{%
  % save equation counter
  \setcounter{alttextimage@saved@equation}{\value{equation}}%
  % (optional) silence labels during the dry run to avoid aux churn
  \let\ATI@orig@label\label
  \let\label\@gobble
  % (optional) if you have a custom \eqlabel, silence it too
  \@ifundefined{eqlabel}{}{%
    \let\ATI@orig@eqlabel\eqlabel
    \zzcommand{\eqlabel}[2]{##2}% just typeset title text if any
  }%
}
\newcommand{\ATI@resume}{%
  % restore equation counter
  \setcounter{equation}{\value{alttextimage@saved@equation}}%
  % restore labeling
  \let\label\ATI@orig@label
  \@ifundefined{ATI@orig@eqlabel}{}{%
    \let\eqlabel\ATI@orig@eqlabel
  }%
}


\newlength{\currentparindent}
\newlength{\currentparskip}
\newlength{\contentheightofalttextimage}
\newlength{\contentheightofalttextimagecaption}
\newsavebox{\alttextbox}

\ExplSyntaxOn
\newdimen\l_tmpa_img_width
\newdimen\l_tmpa_img_height
\newdimen\l_tmpa_ratio
\newdimen\l_tmpa_min_height
\newdimen\l_tmpa_final_height
\ExplSyntaxOff

\ExplSyntaxOn
% \AltTextImageR{item_content}{placeholder.pdf}{caption}{label}{alt_caption}
\zzcommand{\AltTextImageR}[5]{% 
  \ATI@suspend%
  \setcounter{AltTextImageCurrentSide}{1}%
  \setlength{\currentparindent}{\parindent}%
  \setlength{\currentparskip}{\parskip}%
  \sbox{\alttextbox}{\vbox{\hsize=\dimexpr 0.55\linewidth\relax #1}}%
  \setlength{\contentheightofalttextimage}{\ht\alttextbox}%
  \sbox{\alttextbox}{\vbox{\hsize=\dimexpr 0.4\linewidth\relax #3\quad}}%
  \setlength{\contentheightofalttextimagecaption}{\ht\alttextbox}%
  \ATI@resume%
  \noindent
  \begin{minipage}[]{0.55\linewidth}
    \setlength{\parindent}{\currentparindent}
    \setlength{\parskip}{\currentparskip}
    #1
  \end{minipage}%
  \hfill
  % \begin{adjustbox}{width=0.40\linewidth, height=\contentheightofalttextimage, keepaspectratio}
  \begin{minipage}[]{0.40\linewidth} 
      \centering
      \seq_set_split:Nnn \l_tmpa_seq { , } { #2 }%
      \int_set:Nn \l_tmpa_int { \seq_count:N \l_tmpa_seq }%
      \dim_set:Nn \l_tmpa_dim
      {\dim_eval:n{
            \dim_max:nn{.5\contentheightofalttextimage - .5\contentheightofalttextimagecaption}
              {(\contentheightofalttextimage - \contentheightofalttextimagecaption) / \l_tmpa_int}
          }
      }
      % \begin{mdframed}[leftline=true, rightline=false, topline=false, bottomline=false, linecolor=gray]
      \seq_map_inline:Nn \l_tmpa_seq {%
        \includegraphics[width=\linewidth, height=\l_tmpa_dim, keepaspectratio, valign=c]{##1}\\%
      }
      \ifblank{#3}{}{
        \captionsetup{width=\linewidth}
        % \captionsetup{justification=RaggedRight}
        \captionof{figure}[#5]{#3}
      }
      % \end{mdframed}
      \label{#4}
    \end{minipage}%
  % \end{adjustbox}
  \captionsetup{width=\linewidth}
  \smallskip
}
% \AltTextImageL{item_content}{placeholder.pdf}{caption}{label}{alt_caption}
\zzcommand{\AltTextImageL}[5]{%
  \ATI@suspend%
  \setcounter{AltTextImageCurrentSide}{0}
  \setlength{\currentparindent}{\parindent}%
  \setlength{\currentparskip}{\parskip}%
    \sbox{\alttextbox}{\vbox{\hsize=\dimexpr 0.55\linewidth\relax #1}}%
    \setlength{\contentheightofalttextimage}{\ht\alttextbox}%

    \sbox{\alttextbox}{\vbox{\hsize=\dimexpr 0.4\linewidth\relax #3\quad}}%
    \setlength{\contentheightofalttextimagecaption}{\ht\alttextbox}%
  \ATI@resume%
  \noindent
  \begin{minipage}[]{0.40\linewidth}
    % \begin{adjustbox}{max height=\contentheightofalttextimagecaption}
    % \fbox{
      % \strut\vspace*{-\baselineskip}\newline
      \centering
      % Size: \the\contentheightofalttextimagecaption
      % Parse comma-separated images and stack them
      \seq_set_split:Nnn \l_tmpa_seq { , } { #2 }%
      \int_set:Nn \l_tmpa_int { \seq_count:N \l_tmpa_seq }%
      \dim_set:Nn \l_tmpa_dim {
        \dim_max:nn {
          .5 \contentheightofalttextimage  - (\contentheightofalttextimagecaption / 2)
        }{
          (\contentheightofalttextimage / \l_tmpa_int) - (\contentheightofalttextimagecaption / 2)
        }
      }
      % \dim_set:Nn \l_tmpa_dim { (\contentheightofalttextimage / \l_tmpa_int)  - (\contentheightofalttextimagecaption / 2) }%
      % \begin{mdframed}[leftline=false, rightline=false, topline=false, bottomline=false, linecolor=gray]
      \seq_map_inline:Nn \l_tmpa_seq {%
        \includegraphics[width=\linewidth, height=\l_tmpa_dim, keepaspectratio, valign=c]{##1}\\%
      }%
      % \includegraphics[width=\linewidth, height=\contentheightofalttextimage, keepaspectratio, valign=t]{#2}
      \ifblank{#3}{}{
        % \captionsetup{font=footnotesize}
        \captionsetup{width=\linewidth}
        % \captionsetup{justification=RaggedRight}
        \captionof{figure}[#5]{#3}
      }
      % \end{mdframed}
      \label{#4}
    % }
    % \end{adjustbox}
  \end{minipage}%
  \hfill % Add horizontal space between text and image
  \begin{minipage}[]{0.55\linewidth}
    \setlength{\parindent}{\currentparindent}
    \setlength{\parskip}{\currentparskip}
    #1
  \end{minipage}%
  \captionsetup{width=\linewidth}
  \smallskip
}
\ExplSyntaxOff
% \AltTextImage{item_content}{placeholder.pdf}{caption}{label}{alt_caption}
\zzcommand{\AltTextImage}[5]
{
  \stepcounter{AltTextImageCurrentSide}  % Increment the counter
  \ifodd\value{AltTextImageCurrentSide}  % Check if the counter is odd
      \AltTextImageR{#1}{#2}{#3}{#4}{#5}
  \else%
      \AltTextImageL{#1}{#2}{#3}{#4}{#5}
  \fi
}
\ExplSyntaxOff

% \AltTextImage{item_content}{placeholder.pdf}{caption}{label}{alt_caption}
\zzcommand{\AltTextImage}[5]
{
  \stepcounter{AltTextImageCurrentSide}  % Increment the counter
  \ifodd\value{AltTextImageCurrentSide}  % Check if the counter is odd
      \AltTextImageR{#1}{#2}{#3}{#4}{#5}
  \else%
      \AltTextImageL{#1}{#2}{#3}{#4}{#5}
  \fi
}

% \AltTextImageCancelled{item_content}{placeholder.pdf}{caption}{label}{alt_caption}
\zzcommand{\AltTextImageCancelled}[5] {
  \begin{figure}
    \autofitgraphics[]{#2}
    \caption[#5]{#3}
    \label{#4}
  \end{figure}
  #1
}
% \AltTextImageTop{item_content}{placeholder.pdf}{caption}{label}{alt_caption}
\zzcommand{\AltTextImageTop}[5] {
  \begin{figure}%[H]
    \autofitgraphics[]{#2}
    \caption[#5]{#3}
    \label{#4}
  \end{figure}
  #1
}
% \AltTextImageBottom{item_content}{placeholder.pdf}{caption}{label}{alt_caption}
\zzcommand{\AltTextImageBottom}[5] {
  #1
  \begin{figure}%[H]
    \autofitgraphics[]{#2}
    \caption[#5]{#3}
    \label{#4}
  \end{figure}
}

% \lstset {
%     language=C++,
%     backgroundcolor=\color{black!5}, % set backgroundcolor
%     basicstyle=\footnotesize,% basic font setting
% }

\makeglossaries
\DeclareDocumentCommand{\newdualentry}{ O{} O{} m m m m g } {
    \newglossaryentry{gls-#3}{
        name={#5},
        text={#5\if\relax\detokenize{#4}\relax\else\glsadd{#3}\fi},
        plural={\IfValueTF{#7}{#7}{#5s}},
        longplural={\IfValueTF{#7}{#7}{#5s}},
        description={#6},
        #1
    }
    \if\relax\detokenize{#4}\relax
    % nothing
  \else
    \newacronym[longplural={\IfValueTF{#7}{#7}{#5s}}, 
        see={[Glossary:]{gls-#3}},
        #2]{#3}{#4}{#5\glsadd{gls-#3}}
  \fi
    \makeglossaries
}


% \gloss{entry_name}
\zzcommand{\gloss}[1]{\ifglsentryexists{gls-#1}{\glsentryname{gls-#1}\ifcsdef{inheading}{}{{\glsadd{gls-#1}}}}{[XXXX]}}
% \glosses{entry_name}
\zzcommand{\glosses}[1]{\ifglsentryexists{gls-#1}{\glsentrylongpl{gls-#1}\ifcsdef{inheading}{}{{\glsadd{gls-#1}}}}{[XXXX]}}
% \Gloss{entry_name}
\zzcommand{\Gloss}[1]{\ifglsentryexists{gls-#1}{\Glsentryname{gls-#1}\ifcsdef{inheading}{}{{\glsadd{gls-#1}}}}{[XXXX]}}
% \Glosses{entry_name}
\zzcommand{\Glosses}[1]{\ifglsentryexists{gls-#1}{\Glsentrylongpl{gls-#1}\ifcsdef{inheading}{}{{\glsadd{gls-#1}}}}{[XXXX]}}
\zzcommand{\glsnamefont}[1]{\makefirstuc{#1}}


% \crefAnnex{label}
\zzcommand{\crefAnnex}[1]{%
  \begingroup
  \crefname{chapter}{Annex}{Annexes}%
  \Crefname{chapter}{Annex}{Annexes}%
  \cref{#1}%
  \endgroup
}


% \subsubsubsection{title}
\zzcommand{\subsubsubsection}[1]{\bigskip\noindent\textit{#1}\nopagebreak\par}
% \subsubsubsubsection{title}
\zzcommand{\subsubsubsubsection}[1]{\bigskip\noindent\textbf{#1}~}
% % \subsubsubsubsubsection{title}
% \zzcommand{\subsubsubsubsubsection}[1]{\bigskip\noindent\textit{#1}~}

% \wrapFig{img_path}{width}{label}{caption}
\zzcommand{\wrapFig}[4]{
  \wrapFigL{#1}{#2}{#3}{#4}
}

% \wrapFigLR{img_path}{width}{label}{caption}{|L,R|}
\zzcommand{\wrapFigLR}[5]{
  \def\tempwidth{#2}%
  \ifx\tempwidth\empty
    \def\tempwidth{0.25}%
  \fi
  \begin{wrapfigure}{#5}{\tempwidth \linewidth}
    \includegraphics[width=\linewidth]{#1}
    \expandafter\ifx\expandafter\relax\detokenize{#4}\relax\else\caption{#4}\fi
    \expandafter\ifx\expandafter\relax\detokenize{#3}\relax\else\label{#3}\fi
    \vspace{-1 \baselineskip}
  \end{wrapfigure}
}

% \wrapFigL{img_path}{width}{label}{caption}
\zzcommand{\wrapFigL}[4]{
  \wrapFigLR{#1}{#2}{#3}{#4}{L}
}
% \wrapFigR{img_path}{width}{label}{caption}
\zzcommand{\wrapFigR}[4]{
  \wrapFigLR{#1}{#2}{#3}{#4}{R}
}

% \zzcommand{\abstract}
% {
%     \section*{Abstract}
% }

% \shortAbstract{text}
\zzcommand{\shortAbstract}[1]
{
  \noindent
  \textit{#1}
  \bigskip
}

\zzcommand{\bridge}[1]
{
  \subsubsubsection{#1}
}
\zzcommand{\smallConclusion}
{
  \bigskip
  \bigskip
}
  
\zzcommand{\midConclusion}
{
  % \bigskip
  % \noindent\textit{Conclusion}
  % \smallskip
  \subsubsection*{Conclusion}
}
  
\zzcommand{\bigConclusion}
{
  \par\bigskip
  \begin{center}\rule{0.4\linewidth}{0.4pt}\end{center}
  \bigskip\par
}

\ExplSyntaxOn
\NewDocumentCommand{\appendtographicspath}{m}
 {
  \tl_if_exist:cF { Ginput@path } { \tl_new:c { Ginput@path } }
  \tl_gput_right:cn {Ginput@path} { #1 }
 }
\NewDocumentCommand{\prependtographicspath}{m}
 {
  \tl_if_exist:cF { Ginput@path } { \tl_new:c { Ginput@path } }
  \tl_gput_left:cn {Ginput@path} { #1 }
 }
\ExplSyntaxOff

% \resetgraphicspath
\zzcommand{\resetgraphicspath}
{
  \graphicspath{}
}

% \hide{text}
\zzcommand{\hide}[1] {}




% Counter for citations per section
\newcounter{sectioncites}
\newcounter{subsectioncites}
\newcounter{subsubsectioncites}

% Macro to store list of citation keys used in current section
\zzcommand{\sectioncitedkeyslist}{,} % Start with a leading comma
\zzcommand{\subsectioncitedkeyslist}{,} % Start with a leading comma
\zzcommand{\subsubsectioncitedkeyslist}{,} % Start with a leading comma

% Reset at each section
\zzcommand{\resetsectioncites}{%
  \setcounter{sectioncites}{0}%
  \zzcommand{\sectioncitedkeyslist}{,}%
  \resetsubsectioncites
}
\zzcommand{\resetsubsectioncites}{%
  \setcounter{subsectioncites}{0}%
  \zzcommand{\subsectioncitedkeyslist}{,}%
  \resetsubsubsectioncites
}
\zzcommand{\resetsubsubsectioncites}{%
  \setcounter{subsubsectioncites}{0}%
  \zzcommand{\subsubsectioncitedkeyslist}{,}%
}
% Helper: checks if key is in the list, and adds if not
\zzcommand{\adduniquecite}[1]{%
  \def\currentkey{#1}%
  \IfSubStr{\sectioncitedkeyslist}{,\currentkey,}{\def\foundkey{1}}{\def\foundkey{0}}%
  \ifnum\foundkey=0%
    \stepcounter{sectioncites}%
    \xappto{\sectioncitedkeyslist}{,\currentkey,}%
  \fi%
  \IfSubStr{\subsectioncitedkeyslist}{,\currentkey,}{\def\foundkey{1}}{\def\foundkey{0}}%
  \ifnum\foundkey=0%
    \stepcounter{subsectioncites}%
    \xappto{\subsectioncitedkeyslist}{,\currentkey,}%
  \fi%
  \IfSubStr{\subsubsectioncitedkeyslist}{,\currentkey,}{\def\foundkey{1}}{\def\foundkey{0}}%
  \ifnum\foundkey=0%
    \stepcounter{subsubsectioncites}%
    \xappto{\subsubsectioncitedkeyslist}{,\currentkey,}%
  \fi%
}

% Custom bibmacro to do the tracking per entry
% \newbibmacro*{countcite}[1]{%
%   \adduniquecite{#1}%
% }

\zzcommand{\showsectioncites}{%
  \notblank{\currentsection}{%
  \ifdef{\thesectioncites}{%
    \ifnum\thesectioncites=0%
    \else%
      \textcolor{white}{\texttransparent{0.0}{~\thesectioncites~(\currentsection)}}%
    \fi%
  }{}}{}%
}
\zzcommand{\showsubsectioncites}{%
  \notblank{\currentsubsection}{%
  \ifdef{\thesubsectioncites}{%
    \ifnum\thesubsectioncites=0%
    \else%
      \textcolor{white}{\texttransparent{0.0}{~\thesubsectioncites~(\currentsubsection)}}%
    \fi%
  }{}}{}%
}
\zzcommand{\showsubsubsectioncites}{%
  \notblank{\currentsubsubsection}{%
  \ifdef{\thesubsubsectioncites}{%
    \ifnum\thesubsubsectioncites=0%
    \else%
      \textcolor{white}{\texttransparent{0.0}{~\thesubsubsectioncites~(\currentsubsubsection)}}%
    \fi%
  }{}}{}%
}

% \pretocmd{\section}{\showsubsubsectioncites\showsubsectioncites\showsectioncites\resetsectioncites}{}{}
% \pretocmd{\subsection}{\showsubsubsectioncites\showsubsectioncites\resetsubsectioncites}{}{}
% \pretocmd{\subsubsection}{\showsubsubsectioncites\resetsubsubsectioncites}{}{}

\pretocmd{\chapter}{\setcounter{AltTextImageCurrentSide}{0}}{}{}
\pretocmd{\section}{\resetsectioncites}{}{}
\pretocmd{\subsection}{\resetsubsectioncites}{}{}
\pretocmd{\subsubsection}{\resetsubsubsectioncites}{}{}



\zzcommand\theadfont{\bfseries}






% Holds user-defined row width
\newcommand{\fitrowwidth}{\linewidth}
\newcommand{\fitimageoptions}{}

% autofitgraphics with [options]{img1,img2,...}

\NewDocumentCommand{\autofitgraphics}{O{} m}{%
  \zzcommand{\fitrowwidth}{\linewidth}%
  \def\fitimageoptions{}%
  \def\cleanopts{}%
  \def\tempa{#1}%
  \@for\opt:=\tempa\do{%
    \IfBeginWith{\opt}{width=}%
      {\StrBehind{\opt}{=}[{\tempwidth}]\zzcommand{\fitrowwidth}{\tempwidth}}%
      {\ifx\cleanopts\@empty%
         \edef\cleanopts{\opt}%
       \else%
         \edef\cleanopts{\cleanopts,\opt}%
       \fi%
      }%
  }%
  \edef\fitimageoptions{\cleanopts}%
  % \fbox{%
  \fitNimages{#2}%
  % }%
}


\NewDocumentCommand{\labelledautofitgraphics}{O{} m m}{%
  \zzcommand{\fitrowwidth}{\linewidth}%
  \def\fitimageoptions{}%
  \def\cleanopts{}%
  \def\tempa{#1}%
  \@for\opt:=\tempa\do{%
    \IfBeginWith{\opt}{width=}%
      {\StrBehind{\opt}{=}[{\tempwidth}]\zzcommand{\fitrowwidth}{\tempwidth}}%
      {\ifx\cleanopts\@empty%
         \edef\cleanopts{\opt}%
       \else%
         \edef\cleanopts{\cleanopts,\opt}%
       \fi%
      }%
  }%
  \edef\fitimageoptions{\cleanopts}%
  \fitNimagesWithoutPrinting{#3}%
  \addVerticalLabel{#2}%
}

\newcounter{fitImgCount}%
\newcounter{fitImgIndex}%
% --- new: storage for the most-recent row ---
\newsavebox{\fitRowBox}

% --- new: label sizing / styling knobs (tweak as you like) ---
\newlength{\fitLabelWidth}   \setlength{\fitLabelWidth}{1.2em} % column reserved for vertical label
\newlength{\fitLabelPad}     \setlength{\fitLabelPad}{0mm}     % gap between label and images
\newcommand{\fitLabelfont}{\footnotesize}    % label font


\zzcommand{\fitNimages}[1]{%
  \setcounter{fitImgCount}{0}%
  \setcounter{fitImgIndex}{0}%
  \begingroup%
  \def\fitTotalRatio{0}%
  \global\let\fitWidths\empty%
  \def\fitImgList{}%
  % Copy input to a clean list for repeated processing
  \zzcommand{\do}[1]{%
    \edef\fitImgList{\fitImgList,##1}%
  }%
  \docsvlist{#1}%
  % First pass: measure ratios
  \zzcommand{\do}[1]{%
    \stepcounter{fitImgCount}%
    \sbox0{\includegraphics{##1}}%
    \edef\imgWidth{\the\wd0}%
    \edef\imgHeight{\the\ht0}%
    \pgfmathsetmacro{\imgRatio}{\imgWidth/\imgHeight}%
    \pgfmathparse{\fitTotalRatio + \imgRatio}%
    \xdef\fitTotalRatio{\pgfmathresult}%
    \expandafter\xdef\csname fitRatio\thefitImgCount\endcsname{\imgRatio}%
  }%
  \docsvlist{#1}%
  % Second pass: build the row box (width = \fitrowwidth)
  \setcounter{fitImgIndex}{0}%
  % \global\setbox\fitRowBox=\hbox{%
    % \fbox{%
    \begin{adjustbox}{width=\fitrowwidth}%
      \zzcommand{\do}[1]{%
        \stepcounter{fitImgIndex}%
        \pgfmathsetmacro{\relwidth}{\csname fitRatio\thefitImgIndex\endcsname / \fitTotalRatio}%
        % \expandafter\fbox{\expandafter\includegraphics\expandafter[\fitimageoptions, width=\relwidth\linewidth]{##1}}%
        \expandafter\includegraphics\expandafter[\fitimageoptions, width=\relwidth\linewidth]{##1}%
      }%
      \docsvlist{#1}%
    \end{adjustbox}%
    % }%
  % }%
  \endgroup%
}

\zzcommand{\fitNimagesWithoutPrinting}[1]{%
  \setcounter{fitImgCount}{0}%
  \setcounter{fitImgIndex}{0}%
  \begingroup%
  \def\fitTotalRatio{0}%
  \global\let\fitWidths\empty%
  \def\fitImgList{}%
  % Copy input to a clean list for repeated processing
  \zzcommand{\do}[1]{%
    \edef\fitImgList{\fitImgList,##1}%
  }%
  \docsvlist{#1}%
  % First pass: measure ratios
  \zzcommand{\do}[1]{%
    \stepcounter{fitImgCount}%
    \sbox0{\includegraphics{##1}}%
    \edef\imgWidth{\the\wd0}%
    \edef\imgHeight{\the\ht0}%
    \pgfmathsetmacro{\imgRatio}{\imgWidth/\imgHeight}%
    \pgfmathparse{\fitTotalRatio + \imgRatio}%
    \xdef\fitTotalRatio{\pgfmathresult}%
    \expandafter\xdef\csname fitRatio\thefitImgCount\endcsname{\imgRatio}%
  }%
  \docsvlist{#1}%
  % Second pass: build the row box (width = \fitrowwidth)
  \setcounter{fitImgIndex}{0}%
  \global\setbox\fitRowBox=\hbox{%
    \begin{adjustbox}{width=\fitrowwidth}%
      \zzcommand{\do}[1]{%
        \stepcounter{fitImgIndex}%
        \pgfmathsetmacro{\relwidth}{\csname fitRatio\thefitImgIndex\endcsname / \fitTotalRatio}%
        \expandafter\includegraphics\expandafter[\fitimageoptions, width=\relwidth\linewidth]{##1}%
      }%
      \docsvlist{#1}%
    \end{adjustbox}%
  }%
  \endgroup%
}

% --- new: output the cached row with a vertical label on the left ---
\NewDocumentCommand{\addVerticalLabel}{m}{%
  % If someone calls this before \autofitgraphics, we’ll just make an empty box.
  \ifvoid\fitRowBox\relax \global\setbox\fitRowBox=\hbox{} \fi%
  \par\noindent%
  % Scale the combined (label + row) to \linewidth so it never overflows.
  \begin{adjustbox}{width=\fitrowwidth}%
    % \fbox{%
      \begingroup
        \dimen0=\dimexpr\ht\fitRowBox+\dp\fitRowBox\relax
        % Left label column
        % \fbox{
            \vbox to \dimen0{%
              \vfil
              \hbox to \fitLabelWidth{%
                \hfil
                \rotatebox{90}{\centering \fitLabelfont #1}%
                \hfil
              }%
              \vfil
            }%
        % }
        \hspace{\fitLabelPad}%
        % Right: the previously cached row
        \usebox{\fitRowBox}%
      \endgroup
    % }%
  \end{adjustbox}%
  \par\smallskip%
  % Clear the box so the next call starts fresh
  \global\setbox\fitRowBox=\hbox{}%
}

\newcommand{\setfitratios}[1]{%
  \setcounter{fitImgCount}{0}%
  \def\fitTotalRatio{0}%
  \zzcommand{\do}[1]{%
    \stepcounter{fitImgCount}%
    \expandafter\xdef\csname fitRatio\thefitImgCount\endcsname{##1}%
    \pgfmathparse{\fitTotalRatio + ##1}%
    \xdef\fitTotalRatio{\pgfmathresult}%
  }%
  \docsvlist{#1}%
}

\newcounter{fitCapIndex}
% Key-value options for \autofitcaptions
\pgfkeys{
  /autofitcaptions/.is family, /autofitcaptions,
  raise/.store in=\fitraiseamount,
  fontscale/.store in=\fitfontscale,
  % defaults
  raise=2ex,
  fontscale=1,
  % Back-compat: if the user passes a lone token like [1.5ex],
  % treat it as raise=<that>
  .unknown/.code = {\def\fitraiseamount{#1}}
}

\NewDocumentCommand{\autofitcaptions}{O{} m}{%
  % Load defaults, then apply user-specified options (if any)
  \pgfkeys{/autofitcaptions, raise=1em, fontscale=1}%
  \ifblank{#1}{}{%
    \pgfkeys{/autofitcaptions, #1}%
  }%
  \setcounter{fitCapIndex}{0}%
  % \smash[]{%
  \raisebox{\fitraiseamount}{%
  \begin{adjustbox}{width=\fitrowwidth}%
    \zzcommand{\do}[1]{%
      \stepcounter{fitCapIndex}%
      \pgfmathsetmacro{\relwidth}{\csname fitRatio\thefitCapIndex\endcsname / \fitTotalRatio}%
      % \fbox{%
      % \raisebox{\fitraiseamount}{%
        % \parbox[t]{\relwidth\linewidth}{\centering \scalebox{\fitfontscale}{##1}}%
          \scalebox{\fitfontscale}{%
            \parbox[t]{\relwidth\linewidth}{\centering ##1}%
          }%
        % }%
      % }%
      % }%
    }%
    \ifblank{#2}{%
      \docsvlist{~}%
    }{%
      \docsvlist{#2}%
    }%
  \end{adjustbox}%
  }%
  % }%
}


% % \citet{ref}
% \zzcommand{\citet}[1]{\citep{#1}}

% \delete{text_to_remove}
\zzcommand{\delete}[1]
{\remove{#1}}
\zzcommand{\comment}[1]
{\par \highlight{~!~!~!~!~!~!~!~!~!~!~} #1 \highlight{~!~!~!~!~!~!~!~!~!~!~} \\}
\zzcommand{\Replace}{\replace}
\zzcommand{\Comment}{\comment}
\zzcommand{\Add}{\add}
\zzcommand{\Delete}{\delete}
\zzcommand{\revrepl}{\replace}
\zzcommand{\revnote}{\comment}


% \addFileAtTheEnd{tex_file}
\zzcommand{\addFileAtTheEnd}[1]{
  \clearpage
  \thispagestyle{empty}  % désactive l'en-tête et pied
  \newgeometry{top=1.8cm, bottom=1.8cm, left=2cm, right=2cm}

  % Si on est sur une page impaire, on force une page blanche pour aller à une page paire
  \ifodd\value{page}
  \null
  \thispagestyle{empty}
  \newpage
  \fi
  \input{#1}
}


% \newcommand{\listequationsname}{List of Equations}
% \newlistof{myequations}{equ}{\listequationsname}
% % \AtEndEnvironment{equation}{\addcontentsline{equ}{myequations}{Equation~\theequation: NO NAME}}
% % \AtEndEnvironment{align}{\addcontentsline{equ}{myequations}{Equation~\theequation: NO NAME}}


\zzcommand{\eqlabel}[2]{
  \label{#1}%
  \ifblank{#2}{}{\addcontentsline{equ}{myequations}{\cref{#1}: #2}}
}
\makeatother